\section{Introduction}\label{sec:intro}
Photometric redshifts (photo-$z$s) enable studies of galaxy evolution and
cosmology with millions to billions of galaxies. They provide distances
for tracing galaxy populations and mapping the matter distribution with
weak gravitational lensing. Photo-$z$ errors, however, already limit
cosmological analyses. In the Hyper Suprime-Cam (HSC) Year-3 joint
clustering and weak-lensing analysis, accounting for residual photo-$z$
uncertainty lowered the inferred $S_8=\sigma_8\sqrt{\Omega_m/0.3}$ from
$0.791^{+0.018}_{-0.020}$ to $0.763^{+0.040}_{-0.036}$ and roughly doubled
its error bars \citep{miyatake2023,hahn2024}. Current photo-$z$s also exhibit
biases and underestimated uncertainties. Even COSMOS2020, which uses deep
multiband imaging, shows magnitude-dependent biases and apparent
overconfidence for faint galaxies in spectroscopic comparisons
\citep{weaver2022}.

Photo-$z$ calibration will be more demanding for the Vera C.~Rubin
Observatory's Legacy Survey of Space and Time \citep[LSST;][]{ivezic2019}.
LSST will image approximately $18{,}000\,\mathrm{deg}^2$ in six broadband
filters, $ugrizy$, over ten years and is expected to observe approximately
20 billion galaxies. Exploiting this statistical power requires accurate
redshift distributions for the galaxies used in cosmological analyses
\citep{desc2018}. Calibration data that reach the depth and span the galaxy
populations of LSST will be a major bottleneck.

Secure spectroscopic redshifts provide the reference standard for photo-$z$
training and calibration. DESI achieves redshift precisions of approximately
$10$--$40\,\mathrm{km\,s^{-1}}$ for its main galaxy samples \citep{lan2023}.
These errors are negligible compared with broadband photo-$z$ uncertainties,
so reliable spectroscopic measurements can be treated as known labels.
Self-organizing-map methods group galaxies with similar photometry and use
reference redshifts to calibrate their redshift distributions. SOMPZ connects
these groups in deep and wide-field photometry to infer tomographic redshift
distributions \citep{buchs2019,myles2021}. In DES, its reference sample
combines spectroscopy with high-quality multiband photo-$z$s.

The main limitation of spectroscopic calibration is selection. Spectroscopic
samples preferentially contain brighter galaxies and objects with identifiable
spectral features. Photo-$z$ estimators trained on these samples can therefore
inherit their selection biases. Restricting analyses to galaxies represented
by spectroscopy reduces this risk but discards statistical power. The
KiDS-1000 SOM-gold selection, for example, reduced the effective source density
by approximately 20\% \citep{hildebrandt2021,giblin2021}. Obtaining
representative spectroscopy at LSST depths would require substantial telescope
time. \citet{newman2015} estimate that conventional training requires at least
approximately 30,000 spectra across multiple independent fields, reaching
the faint limits of the science sample.

Deep multiband imaging offers an alternative. COSMOS2020 combines more than
30 photometric bands across approximately $2\,\mathrm{deg}^2$, including
12 optical intermediate bands and two narrow bands \citep{weaver2022}.
These measurements constrain redshift more precisely than broadband photometry
alone: normalized redshift scatter is below 1\% for $i<21$ galaxies and
reaches approximately 5\% at $25<i<27$. These photo-$z$s cannot be treated
as exact labels. Their released redshift probability distributions, however,
allow us to propagate their uncertainties into a broadband estimator.
Their calibration must also be tested: \citet{weaver2022} find underestimated
uncertainties at faint magnitudes, while noting that spectroscopic selection
complicates this comparison.

In this work, we present a framework that propagates deep multiband redshift
uncertainties into broadband photo-$z$s using normalizing flows. We sample
redshifts from the deep-field probability distributions and broadband
measurements from a photometric transfer model, then train a conditional flow
on these samples. Normalizing flows can represent complex, multimodal
distributions and provide efficient likelihood evaluation and sampling
\citep{durkan2019,papamakarios2021}. Under a conditional sufficiency assumption,
this procedure marginalizes over the deep observations to estimate
$p(z\given x)$ for broadband measurements $x$. It propagates the input
redshift uncertainties but does not automatically correct their systematic
errors.

We demonstrate this framework with public COSMOS2020 data. We use the released
multiband redshift distributions to train a flow on CFHT $u$/HSC $grizy$
photometry, adding noise where needed to approximate LSST ten-year depths.
The experiment retains the native COSMOS filters and uses an explicit noise
model for the transfer. On a spatially held-out spectroscopic sample, we
compare our predictions with a spectroscopically trained broadband flow,
broadband template fitting, and full-multiband template-fitting estimates.
We test both redshift accuracy and uncertainty calibration to assess whether
probabilistic labels can extend broadband training to larger and fainter
samples.
